\documentclass[12pt]{article}
\usepackage[T1]{fontenc}
\usepackage[utf8]{inputenc}
\usepackage{lmodern}
\usepackage{textcomp}
\usepackage{enumitem}
\usepackage{geometry}
\geometry{margin=1in}
\begin{document}
\begin{center}
Answer Key - Business Administration - Undergraduate Level Assessment 11 \
\small (Answers are ordered exactly as in the question paper.)
\end{center}

\section*{Multiple Choice}
\begin{enumerate}[label=\arabic*.]
\item \textbf{Answer:} C
\\ \textit{Options:} A) Conceptual skills; B) Analytical skills; C) IT and computing skills; D) Communication skills
\\ \textit{Note:} IT and computing skills are considered more technical competencies rather than essential management skills specifically associated with the planning function, which typically emphasizes strategic thinking, resource allocation, and decision-making.
\item \textbf{Answer:} C
\\ \textit{Options:} A) All members should be trained in media relations.; B) A member of senior management should be included in the team.; C) A lawyer should be included in the team.; D) All members should be trained in group decision making.
\\ \textit{Note:} Including a lawyer in a crisis management team is not universally necessary, as the composition of such a team often depends on the nature of the crisis and may prioritize other roles such as communication or operations over legal representation.
\item \textbf{Answer:} A
\\ \textit{Options:} A) 1,2; B) 1,3; C) 2,3; D) 1,2,3
\\ \textit{Note:} The correct answer indicates that the information is utilized for purposes that benefit the public, reflecting the ethical responsibility of businesses to prioritize societal welfare alongside their operational goals.
\item \textbf{Answer:} D
\\ \textit{Options:} A) Capable of compensating for rapid environmental changes and technical improvements.; B) Always available and complete.; C) Seldom obsolete and usually fits the dimensions of your problem.; D) Generally cheaper to gather than primary data and takes less time to find.
\\ \textit{Note:} The correct answer highlights that secondary data is typically more cost-effective and quicker to obtain because it involves using existing data collected by others, rather than the resource-intensive process of gathering original data through primary research methods.
\item \textbf{Answer:} D
\\ \textit{Options:} A) exploratory survey; B) situation interview; C) communication audit; D) stakeholder analysis
\\ \textit{Note:} Stakeholder analysis is the process of identifying and assessing the influence and interests of individuals or groups that are involved with or impacted by an organization's activities, making it the appropriate term for the described process.
\end{enumerate}

\section*{True / False}
\begin{enumerate}[label=\arabic*.]
\setcounter{enumi}{5}
\item \textbf{Answer:} True
\\ \textit{Options:} A) True; B) False
\\ \textit{Note:} The Audit Bureau of Circulation primarily focuses on verifying circulation figures for publications rather than gathering or providing data on readership demographics or behavior.
\item \textbf{Answer:} True
\\ \textit{Options:} A) True; B) False
\\ \textit{Note:} The statement is true because a corporation is a distinct legal entity that exists independently of its shareholders, managers, and employees, enabling it to continue operating indefinitely and to own assets separate from those individuals.
\item \textbf{Answer:} False
\\ \textit{Options:} A) True; B) False
\\ \textit{Note:} Emotional exploitation can significantly harm personal relationships as it often involves manipulation and abuse that undermine trust and emotional well-being, regardless of any financial implications.
\item \textbf{Answer:} False
\\ \textit{Options:} A) True; B) False
\\ \textit{Note:} A crisis management campaign involves addressing and mitigating the effects of a crisis, which often requires navigating differing emotional convictions among stakeholders rather than relying on shared beliefs.
\item \textbf{Answer:} False
\\ \textit{Options:} A) True; B) False
\\ \textit{Note:} The statement is false because while services are indeed characterized by intangibility, variability, and inseparability, they are not inherently characterized by profitability, which can vary based on numerous factors such as market demand and operational efficiency.
\end{enumerate}

\section*{Long Form Response}
\begin{enumerate}[label=\arabic*.]
\setcounter{enumi}{10}
\item \textbf{Answer:} Utilitarianism in business decision-making revolves around the principle of "the greatest good for the greatest number," which suggests that ethical choices should aim to maximize overall happiness and minimize suffering. This approach influences businesses to evaluate the potential outcomes of their decisions not just for their stakeholders, but for the broader community as well. By prioritizing actions that benefit the majority, companies can foster trust and enhance their reputation, thus contributing to long-term success. However, this principle also raises ethical dilemmas, as decisions that favor the majority may inadvertently harm minority groups. Therefore, while utilitarianism provides a framework for making ethical choices, it requires careful consideration of all stakeholders involved.
\\ \textit{Note:} correct because it accurately describes how utilitarianism guides business decision-making by emphasizing the importance of maximizing overall happiness and minimizing suffering for the greatest number of people, which in turn builds trust and enhances a company's reputation.
\item \textbf{Answer:} Businesses can influence government decision-making through various avenues, including lobbying, public relations campaigns, and strategic partnerships. Lobbying involves direct engagement with lawmakers and regulatory agencies, allowing businesses to present their interests and concerns effectively. The breadth of transmission encompasses both formal and informal channels, such as political contributions, grassroots mobilization, and media outreach, which help to amplify their message. The content of communication can range from data-driven arguments and expert testimonies to emotional appeals, all aimed at persuading policymakers to align with business interests. Overall, the combination of these methods enables businesses to shape legislation, regulation, and public policy in ways that support their objectives.
\\ \textit{Note:} correct because it identifies key methods of influence-lobbying, public relations, and partnerships-while discussing the importance of both formal and informal communication channels that businesses use to effectively convey their interests to government decision-makers.
\end{enumerate}

\vfill
\noindent\textit{End of Answer Key}
\end{document}